\chapter{Projekt i implementacja systemu}
\label{chap:system}

System zaprojektowano jako zestaw wymiennych modułów korzystających ze wspólnej reprezentacji danych i wspólnego dekodera klastrów. Dzięki temu wariant bez autokodera, DAE oraz U-Net mogą zostać porównane bez zmiany splitu, scorera i procedury treningowej. Implementację wykonano w Pythonie 3.11 i PyTorch, a parametry eksperymentów przeniesiono do plików YAML. W rozdziale opisano przepływ danych, architektury, funkcje strat, uczenie oraz mechanizmy odtwarzalności.

\section{Architektura potoku}
\label{sec:architektura-potoku}

Wejściem jest dokument po konwersji z CorefUD. Warstwa danych zachowuje słowa, subwordy, globalne identyfikatory wzmianek i maski okien. Enkoder językowy tworzy dla każdej kandydackiej wzmianki wektor $h_i\in\mathbb{R}^{768}$. Reprezentacja łączy stany granic oraz agregację wnętrza, a głowa pozostaje powiązana z indeksem słowa.

Na następnym etapie wykonywany jest jeden z trzech wariantów. Baseline rzutuje $h_i$ bez rekonstrukcji i ocenia pary. DAE koduje oraz rekonstruuje $h_i$, po czym scorer pracuje na kodzie. U-Net nie korzysta bezpośrednio ze scorera wektorowego, lecz segmentuje wielokanałową macierz par. Wszystkie warianty zwracają logity relacji o kształcie $B\times M\times M$.

Maska wskazuje prawdziwe wzmianki, padding i dopuszczalny kierunek pary. Podczas uczenia strata jest liczona na dolnym trójkącie bez diagonali, aby ta sama symetryczna relacja nie występowała dwukrotnie. Podczas dekodowania prawdopodobieństwa są progowane, a union-find tworzy składowe spójne. Próg dobiera się tylko na dev.

Dokument dłuższy od wejścia enkodera jest reprezentowany przez nachodzące okna. Ta sama wzmianka zachowuje globalny identyfikator. Dla pary ocenionej w kilku oknach wybiera się maksimum skalibrowanego prawdopodobieństwa. Dekoder uruchamia się osobno dla każdego dokumentu, co technicznie uniemożliwia utworzenie encji między dokumentami.

Tabela~\ref{tab:tensory-systemu} przedstawia główne tensory. $B$ oznacza batch, $M$ limit wzmianek, $D$ wymiar embeddingu, $F$ liczbę cech par, a $C$ kanały macierzy.

\begin{table}[htbp]
  \centering
  \caption{Główne tensory interfejsu modelu}
  \label{tab:tensory-systemu}
  \begin{tabular}{p{0.30\textwidth}p{0.21\textwidth}p{0.39\textwidth}}
    \toprule
    Tensor & Kształt & Znaczenie \\
    \midrule
    \texttt{mention\_embeddings} & $B\times M\times D$ & kontekstowe reprezentacje wzmianek \\
    \texttt{pair\_features} & $B\times F\times M\times M$ & odległość i zgodność par \\
    \texttt{matrix\_features} & $B\times C\times M\times M$ & wejście segmentacyjne U-Net \\
    \texttt{valid\_mentions} & $B\times M$ & maska prawdziwych pozycji \\
    \texttt{target} & $B\times M\times M$ & złota macierz tożsamości \\
    \texttt{logits} & $B\times M\times M$ & symetryczne oceny modelu \\
    \texttt{latent} & $B\times M\times K$ & opcjonalny kod wzmianki \\
    \texttt{reconstruction} & $B\times M\times D$ & opcjonalna rekonstrukcja DAE \\
    \bottomrule
  \end{tabular}
\end{table}

\section{Wariant kontrolny bez autokodera}
\label{sec:system-baseline}

Kontrola neuronowa rzutuje embedding $D$-wymiarowy do przestrzeni $K$ przez warstwę liniową, normalizację warstwową i GELU. Nie rekonstruuje wejścia. Dla pary $(i,j)$ budowany jest wektor
\[
q_{ij}=[z_i;z_j;z_i\odot z_j;|z_i-z_j|;f_{ij}],
\]
gdzie $f_{ij}$ zawiera jawne cechy relacyjne. Dwuwarstwowa sieć kończy się jednym logitem. Macierz jest symetryzowana przez średnią logitu i jego transpozycji.

Wariant ma taki sam wymiar latentny i taką samą głowicę jak DAE. Różni się wyłącznie brakiem dekodera, szumu i celu rekonstrukcyjnego. Jest zatem właściwą kontrolą PB1, w odróżnieniu od porównania z większym modelem o innym enkoderze.

Symetryzacja jest wykonywana przed stratą. Nie narzuca przechodniości; ta powstaje dopiero w dekoderze. Rozdzielenie pozwala policzyć naruszenia w surowej macierzy oraz ocenić, czy kara globalna wnosi korzyść.

\section{Odszumiający autokoder wzmianek}
\label{sec:system-dae}

DAE ma układ 768--384--128--384--768. Koder zawiera warstwę liniową, LayerNorm, GELU, dropout, drugą projekcję i normalizację kodu. Dekoder odtwarza wymiar wejścia przez analogiczną ścieżkę. W czasie ewaluacji zakłócenie jest wyłączone.

Podczas treningu niezależnie maskuje się 15\% wymiarów oraz dodaje szum Gaussa o odchyleniu 0,01. Ustawienia są parametrami YAML, a nie stałymi ukrytymi w pętli. Plan domenowy przewiduje również maskowanie części cech morfosyntaktycznych, gdy staną się osobnym segmentem wejścia.

Rekonstrukcja łączy błąd $L_1$ i stratę cosinusową na prawdziwych wzmiankach:
\[
\mathcal{L}_{rec}
=
\lVert h-\hat h\rVert_1
+1-\cos(h,\hat h).
\]
Wspólny cel ma postać
\[
\mathcal{L}_{C}
=\mathcal{L}_{edge}
+\lambda_{rec}\mathcal{L}_{rec}.
\]
W implementacji komponent cosinusowy jest częścią funkcji rekonstrukcji, a waga z YAML steruje całością. Padding nie uczestniczy w obliczeniu.

Scorer otrzymuje kod $z$ o wymiarze 128 oraz identyczne cechy par co kontrola. Pretraining i fine-tuning mogą być rozdzielone checkpointem. W obecnym potoku syntetycznym wykonywany jest wspólny trening techniczny; nie zastępuje on planowanego domenowego pretrainingu na prawdziwych embeddingach.

\section{Segmentacja macierzy przez U-Net}
\label{sec:system-unet}

Macierz wejściowa ma domyślnie osiem kanałów. Obejmuje podobieństwo embeddingów, zgodność rodzaju, liczby i osoby, odległość tokenową i zdaniową, zgodność głowy lub lematu oraz oceny pojedynczych wzmianek. W implementacji syntetycznej liczba kanałów jest konfigurowalna, dzięki czemu test nie wymaga prawdziwego taggera.

Kompaktowy U-Net zawiera dwa poziomy zmniejszania rozdzielczości i most. Blok składa się z dwóch splotów $3\times3$, GroupNorm i GELU. Max-pooling zmniejsza rozdzielczość, a splot transponowany ją odtwarza. Mapy z enkodera są konkatenowane z odpowiadającymi poziomami dekodera. Ostatni splot $1\times1$ zwraca jeden kanał.

Jeżeli $M$ nie jest podzielne przez cztery, wejście jest tymczasowo dopełniane. Wynik zostaje przycięty do pierwotnego rozmiaru i symetryzowany:
\[
Z\leftarrow\frac{Z+Z^\top}{2}.
\]
Dopełnienie jest wyłączone z maski straty.

Funkcja celu ma trzy składniki:
\[
\mathcal{L}_{A}
=
\mathcal{L}^{w}_{BCE}
+\lambda_{dice}\mathcal{L}_{Dice}
+\lambda_{tri}\mathcal{L}_{trans}.
\]
Ważona BCE zwiększa koszt błędu dodatnich, rzadkich krawędzi. Dice mierzy miękkie pokrycie dodatniej maski. Kara przechodniości wykorzystuje prawdopodobieństwa trójek i penalizuje dwie silne krawędzie bez trzeciej. Obliczenie jest sześcienne względem $M$, dlatego limit kandydatów i ablacja $\lambda_{tri}=0$ są ważne dla profilowania.

Projektowy U-Net ma kanały 32--64--128 i mniej niż milion parametrów samego modułu. Szacunki 6--10 GB dla zamrożonego enkodera oraz 14--20 GB przy wspólnym strojeniu nie są pomiarami. Kod zapisuje rzeczywistą wartość \texttt{torch.cuda.max\_memory\_allocated}; dopiero ona może trafić do wyników.

\section{Selektywna hybryda}
\label{sec:system-hybryda}

Hybryda działa nad logitami DAE. Dla anafory sortuje wcześniejszych kandydatów i liczy margines dwóch najlepszych ocen. Mały margines oznacza niepewność decyzji. Drugim sygnałem jest błąd rekonstrukcji wzmianki, interpretowany jako odległość od rozkładu poznanego przez DAE.

Klasa \texttt{HybridSelector} wybiera przypadki poniżej progu marginesu lub powyżej opcjonalnego progu błędu rekonstrukcji, porządkuje je rosnąco według marginesu i malejąco według błędu, a następnie ogranicza listę do maksymalnej liczby kandydatów. Domyślny budżet projektu wynosi 16 na dokument. Sam selektor nie wykonuje wywołania sieciowego, co oddziela politykę wyboru od dostawcy LLM.

Adapter API wysyła instrukcję, tekst po audycie oraz listę identyfikatorów wzmianek. W zero-shot lista demonstracji jest pusta, a w few-shot zawiera przykłady wyłącznie z train. Odpowiedź ma format JSON z identyfikatorem anafory oraz poprzednika albo \texttt{null}. Brak klucza środowiskowego, niepoprawny schemat lub obcy identyfikator przerywa decyzję zewnętrzną.

Warstwa LLM nie ma trenowanych parametrów w lokalnym modelu. Koszt mierzy się wywołaniami, tokenami i opóźnieniem. Tekst z potencjalnym PII nie jest wysyłany. Jeśli kontrola nie przejdzie, zachowuje się predykcję lokalną.

\section{Uczenie}
\label{sec:uczenie-systemu}

Pętla treningowa tworzy dataset z pliku tensorowego albo generatora syntetycznego. DataLoader używa osobnego generatora z seedem. AdamW otrzymuje learning rate i weight decay z YAML. Obsługiwane są akumulacja gradientu, clipping normy, mixed precision i automatyczny wybór CPU/CUDA.

Seedy są ustawiane dla Pythona, NumPy i PyTorch. Deterministyczny tryb wyłącza benchmark cuDNN oraz ustawia przestrzeń roboczą CuBLAS. Każdy krok zapisuje epokę, numer i stratę do JSONL. Po treningu powstaje checkpoint ze stanem modelu, konfiguracją i seedem oraz podsumowanie z liczbą parametrów, końcową stratą, urządzeniem i szczytem pamięci.

Generator syntetyczny tworzy prototypy klastrów, embeddingi z szumem i złotą macierz. Służy tylko do testów kształtu i przepływu gradientu. Flaga \texttt{synthetic\_data} trafia do raportu, aby techniczna liczba nie została pomylona z wynikiem badawczym.

\section{Konfiguracja, testy i odtwarzalność}
\label{sec:odtwarzalnosc-systemu}

Pliki \texttt{baseline.yaml}, \texttt{dae.yaml} i \texttt{matrix\_unet.yaml} zawierają wymiary, dane, parametry treningu i wagi strat. \texttt{smoke.yaml} zmniejsza wymiar do 32 i uruchamia 20 przykładów na CPU. Konfiguracje pełne wskazują plik tensorów zamiast cichego tworzenia danych.

Testy sprawdzają kształty wszystkich wariantów, symetrię wyjścia, limit selektora i rzeczywisty krok treningu. Osobny zestaw kontroluje metryki, bootstrap, parser oficjalnego scorera i baseline’y. Pinned scorer jest pobierany przez skrypt i weryfikowany pełnym hashem commitu.

Tabela~\ref{tab:artefakty-systemu} wskazuje główne artefakty reprodukcji.

\begin{table}[htbp]
  \centering
  \caption{Artefakty implementacji i ich funkcje}
  \label{tab:artefakty-systemu}
  \begin{tabular}{p{0.39\textwidth}p{0.51\textwidth}}
    \toprule
    Artefakt & Funkcja \\
    \midrule
    \path{src/models/coreference.py} & DAE, scorer par, U-Net i selektor \\
    \texttt{train.py} & trening, straty, seedy, logi i checkpoint \\
    \texttt{eval.py} & ewaluacja sterowana YAML \\
    \path{src/eval/official.py} & izolowane wywołanie CorefUD scorer \\
    \path{src/eval/istotnosc.py} & bootstrap dokumentowy i porównanie sparowane \\
    \path{src/baselines} & reguła, regresja, Stanza i LLM \\
    \path{scripts/run_reduced_experiments.py} & zredukowane E1--E5 \\
    \bottomrule
  \end{tabular}
\end{table}

\section{Podsumowanie}

Zaimplementowano wspólny interfejs dla kontroli, DAE i U-Net, symetryczne logity oraz dokumentowy dekoder klastrów. DAE uczy rekonstrukcji zaburzonych embeddingów, U-Net segmentuje macierz relacji, a selektor ogranicza użycie LLM. Konfiguracje, seedy, logi, checkpointy i przypięty scorer tworzą podstawę odtwarzalności. Następny rozdział opisuje protokół, który oddziela techniczną weryfikację od eksperymentów rozstrzygających.
